\section{Rectificador de Media Onda}

\subsection{Gráficos de Canales Sincronizados}
Se incluyen a continuación los gráficos obtenidos al medir simultáneamente la señal de entrada (canal 1) y la señal de salida (canal 2) del rectificador de media onda para cada uno de los casos estudiados: sin capacitor, con capacitor de 100 $\mu$F, con capacitor de 220 $\mu$F y con capacitor de 2200 $\mu$F.\\
Los gráficos únicamente corresponden a los experimentos realizados con la resistencia de prueba, ya que las formas de onda obtenidas con el cooler fueron similares, aunque con una caída de tensión ligeramente mayor debido a la naturaleza inductiva del motor.

% --- GRÁFICO 1: SIN CAPACITOR ---
\begin{figure}[htbp]
    \centering
    \begin{tikzpicture}
        \begin{axis}[
            width=0.85\textwidth,
            height=6cm,
            grid=both,
            grid style={dashed, gray!30},
            xlabel={Tiempo (ms)},
            ylabel={Tensión (V)},
            title={Medición de Canales Sincronizados - Sin Filtrar},
            legend pos=north east,
            xmin=0
        ]
            \addplot[color=red, thick] table [x=Tiempo, y=Ch1, col sep=comma] {./csvs/default1.csv};
            \addplot[color=teal, thick] table [x=Tiempo, y=Ch2, col sep=comma] {./csvs/default1.csv};
            \legend{Entrada, Salida}
        \end{axis}
    \end{tikzpicture}
    \caption{Señales de entrada y salida del rectificador de media onda sin capacitor.}
    \label{fig:media_onda_sin_capacitor}
\end{figure}

% --- GRÁFICO 2: CAPACITOR 100uF ---
\begin{figure}[htbp]
    \centering
    \begin{tikzpicture}
        \begin{axis}[
            width=0.85\textwidth,
            height=6cm,
            grid=both,
            grid style={dashed, gray!30},
            xlabel={Tiempo (ms)},
            ylabel={Tensión (V)},
            title={Medición de Canales Sincronizados - C = 100 $\mu$F},
            legend pos=north east,
            xmin=0
        ]
            \addplot[color=red, thick] table [x=Tiempo, y=Ch1, col sep=comma] {./csvs/default2.csv};
            \addplot[color=teal, thick] table [x=Tiempo, y=Ch2, col sep=comma] {./csvs/default2.csv};
            \legend{Entrada, Salida}
        \end{axis}
    \end{tikzpicture}
    \caption{Señales de entrada y salida del rectificador de media onda con capacitor de 100 $\mu$F.}
    \label{fig:media_onda_capacitor_100u}
\end{figure}

% --- GRÁFICO 3: CAPACITOR 220uF ---
\begin{figure}[htbp]
    \centering
    \begin{tikzpicture}
        \begin{axis}[
            width=0.85\textwidth,
            height=6cm,
            grid=both,
            grid style={dashed, gray!30},
            xlabel={Tiempo (ms)},
            ylabel={Tensión (V)},
            title={Medición de Canales Sincronizados - C = 220 $\mu$F},
            legend pos=north east,
            xmin=0
        ]
            \addplot[color=red, thick] table [x=Tiempo, y=Ch1, col sep=comma] {./csvs/default5.csv};
            \addplot[color=teal, thick] table [x=Tiempo, y=Ch2, col sep=comma] {./csvs/default5.csv};
            \legend{Entrada, Salida}
        \end{axis}
    \end{tikzpicture}
    \caption{Señales de entrada y salida del rectificador de media onda con capacitor de 220 $\mu$F.}
    \label{fig:media_onda_capacitor_220u}
\end{figure}

% --- GRÁFICO 4: CAPACITOR 2200uF ---
\begin{figure}[htbp]
    \centering
    \begin{tikzpicture}
        \begin{axis}[
            width=0.85\textwidth,
            height=6cm,
            grid=both,
            grid style={dashed, gray!30},
            xlabel={Tiempo (ms)},
            ylabel={Tensión (V)},
            title={Medición de Canales Sincronizados - C = 2200 $\mu$F},
            legend pos=north east,
            xmin=0
        ]
            \addplot[color=red, thick] table [x=Tiempo, y=Ch1, col sep=comma] {./csvs/default6.csv};
            \addplot[color=teal, thick] table [x=Tiempo, y=Ch2, col sep=comma] {./csvs/default6.csv};
            \legend{Entrada, Salida}
        \end{axis}
    \end{tikzpicture}
    \caption{Señales de entrada y salida del rectificador de media onda con capacitor de 2200 $\mu$F.}
    \label{fig:media_onda_capacitor_2200u}
\end{figure}

\clearpage % Fuerza a que los gráficos se acomoden antes del cuestionario

\subsection{Cuestionario y Análisis de Resultados}

\begin{enumerate}
    \item \textbf{Cálculo del Valor Eficaz ($V_{RMS}$):}
    Para el valor de amplitud cero a pico ($v_{op}$) obtenido sin colocar el capacitor de filtrado, obtenga analíticamente el valor eficaz de la tensión de salida y compare con el medido por medio del multímetro digital. El valor RMS teórico se calcula como:
    
    \[ V_{RMS} = \sqrt{\frac{1}{T} \int_{0}^{T/2} v_{op}^2 \sin^2(\omega t) \, dt} \]
    
    siendo $T$ el período de la señal (considere $f = 50\text{ Hz} \implies T = 20\text{ ms}$).
    
    \begin{tcolorbox}[colback=white, colframe=gray!50, title=Respuesta]
        Desarrollo de la integral: $V_{RMS} = \frac{v_{op}}{2}$ \\
        Valor teórico calculado (Gráfico \ref{fig:media_onda_sin_capacitor}): $V_{RMS} = \text{7.40 V}$. \\
        Valor medido con multímetro: 5.00 V para el caso de la resistencia y 5.15 V para el caso del cooler. \\
        
        La diferencia entre el valor analítico eficaz ($7.40\text{ V}$) y las mediciones del multímetro no se debe a un error instrumental, sino a que se están comparando dos parámetros físicos distintos. Al realizar la medición en la escala de tensión continua (DC), el multímetro registra el valor medio ($\bar{v}$) de la señal y no su valor eficaz ($V_{RMS}$). 
        
        Si se compara la medición con el valor medio teórico calculado en el punto 6 ($\bar{v} = 4.71\text{ V}$), se observa una coincidencia muy estrecha. La pequeña discrepancia restante se atribuye a la caída de tensión en el diodo real y a las tolerancias de los componentes.\\
        Por otra parte, el valor ligeramente mayor en el cooler ($5.15\text{ V}$) responde a su comportamiento como carga inductiva activa, cuya fuerza contraelectromotriz altera levemente la forma de la onda pulsante respecto a la resistencia pura.
    \end{tcolorbox}
    \item \textbf{Relación entre Rizado y Capacitancia:}
    ¿Encuentra alguna relación entre la amplitud del rizado y la capacitancia del capacitor?
    
    \begin{tcolorbox}[colback=white, colframe=gray!50, title=Respuesta]
        A partir de los gráficos \ref{fig:media_onda_capacitor_100u}, \ref{fig:media_onda_capacitor_220u} y \ref{fig:media_onda_capacitor_2200u} se observa que el rizado disminuye a medida que aumenta la capacitancia. Esto se debe a que el capacitor actúa como un reservorio de carga, suavizando las caídas de tensión entre los picos de la señal rectificada. Matemáticamente, el rizado es inversamente proporcional a la capacitancia, lo que explica esta relación.
    \end{tcolorbox}

    \item \textbf{Capacitor de Filtrado:}
    ¿Por qué se denomina ``capacitor de filtrado''? ¿Qué filtra?
    
    \begin{tcolorbox}[colback=white, colframe=gray!50, title=Respuesta]
        Se denomina ``capacitor de filtrado'' porque actúa como un filtro pasa-bajos que atenúa las variaciones bruscas de tensión (componentes de alterna o rizado) provenientes de la rectificación. El capacitor funciona almacenando energía eléctrica durante los intervalos de conducción de los diodos (cuando la tensión sube) y entregándola gradualmente a la carga cuando la tensión de la fuente decrece. Esto suaviza las fluctuaciones y estabiliza la salida transformando los pulsos en una tensión continua lo más constante posible.
    \end{tcolorbox}

    \item \textbf{Velocidad de Rotación del Cooler:}
    ¿Por qué aumenta la velocidad de rotación del cooler al aumentar la capacitancia del capacitor?
    
    \begin{tcolorbox}[colback=white, colframe=gray!50, title=Respuesta]
        Al aumentar la capacitancia, la tensión de la señal no cae abruptamente a cero, esto genera que el cooler reciba una tensión más constante y cercana a su valor máximo durante un período más largo. Esto se traduce en una mayor potencia entregada al cooler, lo que a su vez hace que su velocidad de rotación aumente.
    \end{tcolorbox}

    \item \textbf{Relación con la Medición del Multímetro:}
    Relacione el punto anterior con la tensión de salida medida con el multímetro en la escala de continua.
    
    \begin{tcolorbox}[colback=white, colframe=gray!50, title=Respuesta]
        Al aumentar la capacitancia los valores obtenidos con el multímetro fueron:
        \begin{itemize}
            \item Sin capacitor: $\text{5.00 V}$ (resistencia) y $\text{5.15 V}$ (cooler).
            \item Con capacitor de 100 $\mu$F: $\text{8.87 V}$ (resistencia) y $\text{7.69 V}$ (cooler).
            \item Con capacitor de 220 $\mu$F: $\text{10.90 V}$ (resistencia) y $\text{9.76 V}$ (cooler).
            \item Con capacitor de 2200 $\mu$F: $\text{12.40 V}$ (resistencia) y $\text{11.60 V}$ (cooler).
        \end{itemize}
        Se observa que el valor medio de la tensión de salida aumenta a medida que se incrementa la capacitancia, lo que se traduce en una mayor potencia entregada al cooler y, por ende, un aumento en su velocidad de rotación.
    \end{tcolorbox}

    \item \textbf{Factor de Rizado y Eficiencia Energética:}
    La señal de salida puede escribirse como:
    \[ v(t) = \bar{v} + v_{\text{rizado}}(t) \]
    siendo:
    \[ \bar{v} = \frac{1}{T} \int_{0}^{T/2} v_{op} \sin(\omega t) \, dt \]
    el valor medio de la señal de salida, con $v_{op}$ la tensión 0 a pico de la señal de media onda y $T$ su período. Es decir, la señal de salida puede interpretarse como una señal continua constante con un rizado dependiente del tiempo superpuesto.
    
    Para la salida obtenida sin el capacitor de filtrado, calcule analíticamente el factor de rizado $r$, definido como:
    \[ r = \frac{v_{\text{rizado}, RMS}}{\bar{v}} \]
    ¿Qué puede concluir con respecto a la eficiencia energética del rectificador de media onda?
    
    \begin{tcolorbox}[colback=white, colframe=gray!50, title=Respuesta]
        El valor medio calculado es $\bar{v} = \frac{v_{op}}{\pi}$. \\
        Utilizando el valor eficaz calculado en el punto 1, el valor eficaz del rizado se obtiene a partir de la relación $v_{\text{rizado}, RMS} = \sqrt{V_{RMS}^2 - \bar{v}^2}$.\\\\
        De esta forma se obtuvo: $\bar{v} = \text{4.71 V}$ y $v_{\text{rizado}, RMS} = \text{5.71 V}$. \\
        \\
        Así, calculamos el factor de rizado: $r = \frac{v_{\text{rizado}, RMS}}{\bar{v}} \approx 1.21$ (121\%). \\
        Con respecto a la eficiencia energética, se concluye que el rectificador de media onda sin filtrado tiene una eficiencia muy baja debido a la alta proporción de rizado en comparación con el valor medio de la señal de salida. Esto implica que una gran parte de la energía entregada a la carga se desperdicia en forma de fluctuaciones, lo que no es ideal para aplicaciones que requieren una tensión constante y estable.
    \end{tcolorbox}
\end{enumerate}